\section{Critical-cutoff ledger and route decisions}
\label{sec:tpc61-endpoint-ledger}

The exposure and Gram identities are exact at every finite \(X\).  Their
role in the endpoint argument is determined only after missing mass,
restricted grouping, and the represented reference have been measured on
the same carrier and with the same quantifier.  This section performs that
comparison and records which numerical differences are genuine tests rather
than additive losses.

\subsection{The critical arithmetic wedge}

At the smallest very-high cutoff allowed by the fixed support constants, put
\begin{equation}
 \theta_J:=\frac{133}{400},
 \qquad
 \lambda_L:=\frac{99979}{210000}.
 \label{eq:tpc61-ledger-basic-exponents}
\end{equation}
Thus \(J=X^{\theta_J+o(1)}\), and a critical cofactor block has
\(R=X^{\theta_J+o(1)}\).  The coefficient-blind residual determinant-ladder
degree on that block has exponent
\begin{equation}
 \mu_{\rm res}^{\rm amb}
 :=\lambda_L-\theta_J
 =\frac{15077}{105000}
 =0.143590\ldots .
 \label{eq:tpc61-ledger-critical-residual-degree}
\end{equation}
This number is much larger than \(1/400\), but it is not by itself an
endpoint loss.  In the mass--degree certificate it is multiplied by the
missing mass before the reverse triangle inequality is applied.

Suppose the missing set in the critical block has cardinality
\begin{equation}
 |E_m(R)|\le X^{e+o(1)}
 \label{eq:tpc61-ledger-critical-missing-count}
\end{equation}
uniformly in the stated row class, and suppose the balanced-exposure
hypothesis \(H_m^{\rm eff}\ge X^{\theta_J+o(1)}\) holds there.  Then
\cref{cor:tpc61-dyadic-saving-exponent} gives the conditional missing-mass
saving
\begin{equation}
 \sigma_{\rm cof}=\theta_J-e.
 \label{eq:tpc61-ledger-critical-missing-saving}
\end{equation}
Consequently the residual degree and cofactor saving leave the algebraic
window
\begin{equation}
 \boxed{
 2\theta_J-\lambda_L
 =\frac{39671}{210000}
 =0.188909\ldots .}
 \label{eq:tpc61-ledger-residual-window}
\end{equation}

Let the actual residual-label erasure on this block satisfy
\begin{equation}
 \beta_{M,R}\le X^{\mu_{\rm erase}+o(1)},
 \label{eq:tpc61-ledger-erasure-input}
\end{equation}
and let the represented map obey
\begin{equation}
 \|A_X\zeta\|^2
 \ge X^{-\lambda_{\rm rep}+o(1)}\cD_H(\zeta)
 \label{eq:tpc61-ledger-represented-input}
\end{equation}
in the same uniform or distinguished scope.  Paying the ambient residual
degree, but only the restricted erasure norm, gives
\begin{equation}
 \frac{\|M_{X,R}\zeta\|^2}{\|A_X\zeta\|^2}
 \le
 X^{e+\mu_{\rm erase}+\lambda_{\rm rep}
      -(2\theta_J-\lambda_L)+o(1)}.
 \label{eq:tpc61-ledger-residual-product}
\end{equation}
Hence this certificate has zero fixed-power reassembly cost under the strict
condition
\begin{equation}
 \boxed{
 e+\mu_{\rm erase}+\lambda_{\rm rep}
 <\frac{39671}{210000}.}
 \label{eq:tpc61-ledger-residual-door}
\end{equation}
For \(O(\log X)\) compatible blocks, the same conclusion follows only when
the corresponding exponent inequality holds block by block with one fixed
positive margin \(\delta>0\).  Then each block norm is
\(O(X^{-\delta/2+o(1)}\|A_X\zeta\|)\), and the triangle inequality over
the logarithmic family remains \(o(\|A_X\zeta\|)\).  Mere pointwise
strictness with a margin tending to zero is not a uniform asymptotic
certificate.  None of the three hypotheses in
\eqref{eq:tpc61-ledger-critical-missing-count}--%
\eqref{eq:tpc61-ledger-represented-input} is proved by the numerical window
itself.

\subsection{Why the direct ambient door remains closed}

The inherited coefficient-blind native degree has exponent \(267/400\).
Even the smallest nonempty missing set at the critical cofactor scale has
\(e=0\), so the balanced harmonic certificate saves at most
\(\theta_J=133/400\).  A direct phase-blind payment therefore leaves
\begin{equation}
 \boxed{
 \frac{267}{400}-\frac{133}{400}
 =\frac{67}{200}>0.}
 \label{eq:tpc61-ledger-direct-native-gap}
\end{equation}
This proves that the particular proof architecture which combines the best
pure-cardinality missing saving with the full ambient native degree cannot
make the perturbation \(o(1)\).  It does not lower-bound \(\kappa_M\), does
not rule out cancellation in the parity kernel, and is not a loss to be
entered into the final ledger.

The residual window in \eqref{eq:tpc61-ledger-residual-window} and the direct
gap in \eqref{eq:tpc61-ledger-direct-native-gap} therefore have different
logical meanings.  The first is conditional room for a two-stage proof; the
second closes one support-only proof.  Neither is an arithmetic lower bound
for the full physical map.

\subsection{The optimal one-stage effective-product test}

The restricted Gram eliminates the need to commit in advance to the
residual factorization.  Fix a reference diagonal \(\cD_{\rm ref}\) and
suppose, in one declared scope,
\begin{align}
 \cD_M(\zeta)
 &\le X^{-\sigma_M+o(1)}\cD_{\rm ref}(\zeta),
 \label{eq:tpc61-ledger-missing-input}\\
 \kappa_M
 &\le X^{\nu_M+o(1)},
 \label{eq:tpc61-ledger-kappa-input}\\
 \|A_X\zeta\|^2
 &\ge X^{-\lambda_{\rm rep}+o(1)}
        \cD_{\rm ref}(\zeta).
 \label{eq:tpc61-ledger-reference-input}
\end{align}
Then the exact perturbation parameter is bounded by
\begin{equation}
 q_X(\zeta)
 :=\frac{\kappa_M\cD_M(\zeta)}{\|A_X\zeta\|^2}
 \le X^{\nu_M+\lambda_{\rm rep}-\sigma_M+o(1)}.
 \label{eq:tpc61-ledger-q-bound}
\end{equation}
Thus
\begin{equation}
 \boxed{
 \sigma_M>\nu_M+\lambda_{\rm rep}
 \quad\Longrightarrow\quad
 \lambda_{\rm reasm}=0}
 \label{eq:tpc61-ledger-zero-cost-test}
\end{equation}
for a coefficient-uniform statement, and the analogous conclusion holds
with \(\kappa_M(\zeta_*)\) and \(\lambda_{\rm reasm,*}=0\) for one
distinguished physical vector.  These scopes cannot be interchanged.

At equality of exponents, one must use the actual quantity \(q_X\).  If
\begin{equation}
 1-\sqrt{q_X(\zeta)}\ge X^{-a+o(1)},
 \label{eq:tpc61-ledger-critical-margin}
\end{equation}
then the reverse-triangle certificate charges at most
\begin{equation}
 \lambda_{\rm reasm}\le2a
 \label{eq:tpc61-ledger-critical-cost}
\end{equation}
in the matching scope.  If \(q_X\ge1\), that certificate is silent.  It
does not imply destructive interference, because the exact cross term may
still be favorable and the TPC--60 shorted-Gram route remains available.

The rank obstruction is a useful early test.  If
\(r_M=X^{\rho_M+o(1)}\) and
\(\rank F_M\le X^{\omega_M+o(1)}\), then
\begin{equation}
 \nu_M\ge(\rho_M-\omega_M)_+.
 \label{eq:tpc61-ledger-rank-floor}
\end{equation}
Accordingly, the uniform mass--degree route cannot satisfy
\eqref{eq:tpc61-ledger-zero-cost-test} unless
\begin{equation}
 \sigma_M>(\rho_M-\omega_M)_++\lambda_{\rm rep}.
 \label{eq:tpc61-ledger-rank-necessary}
\end{equation}
Failure closes this certificate but need not close a distinguished vector
orthogonal to the top singular directions.

\subsection{The complete endpoint ledger}

The prospective fixed-shift transfer chain remains the one inherited from
TPC--60.  Its losses are recorded as
\begin{equation}
 \boxed{
 \begin{aligned}
 \Lambda_{\rm tot}:={}&
 \lambda_{\rm sel}+\lambda_{\rm occ}
 +\lambda_{\rm blk}+\lambda_{\rm term}
 +\lambda_{\rm sing}\\
 &+\lambda_H+\lambda_{\rm aff}
 +\lambda_{\rm reasm}
 +\lambda_{\rm bdry}+\lambda_{\rm tail}.
 \end{aligned}}
 \label{eq:tpc61-complete-endpoint-ledger}
\end{equation}
The endpoint requirement is
\begin{equation}
 \boxed{\Lambda_{\rm tot}<\frac1{400}.}
 \label{eq:tpc61-strict-endpoint-budget}
\end{equation}
For a distinguished chain, \(\lambda_{\rm reasm}\) is consistently
replaced by \(\lambda_{\rm reasm,*}\).  Equality is rejected unless the
leading constants and all remainders are controlled strongly enough to
recover a strict final inequality.

In particular, in the equality-scale situation
\eqref{eq:tpc61-ledger-critical-margin}--%
\eqref{eq:tpc61-ledger-critical-cost}, put
\(\Lambda_{\rm other}:=\Lambda_{\rm tot}-\lambda_{\rm reasm}\).  The
reverse-triangle certificate can fit the endpoint only if
\begin{equation}
 \boxed{
  2a+\Lambda_{\rm other}<\frac1{400}.}
 \label{eq:tpc61-equality-margin-endpoint-test}
\end{equation}
Thus the number \(a\) is not free slack: it consumes twice its value in
the loss ledger.

The quantities \(e\), \(\sigma_M\), \(\nu_M\),
\(\mu_{\rm erase}\), and \(67/200\) are not additional arrows in
\eqref{eq:tpc61-complete-endpoint-ledger}.  They are inputs or diagnostics
used to decide the single reassembly loss.  When \(\lambda_{\rm rep}\) is
assembled from \(\lambda_H\), \(\lambda_{\rm aff}\), or other arrows already
listed in that ledger, it enters the comparison tests above but is not charged
a second time.  Likewise, the cutoff
\(y\asymp Q\), exact coordinate changes, Moore--Penrose restriction to the
active support, and finite or logarithmic block decompositions have zero
fixed-power cost.

\subsection{Stop and pivot rules}

The current proof architecture has the following exact decisions.
\begin{enumerate}[leftmargin=2.3em]
 \item If some active row has \(W_R(m)=0<W_M(m)\), the
       coefficient-uniform represented-mass route is dead.  An average over
       other rows cannot repair it.
 \item If the dimension exponents force
       \[
        (\rho_M-\omega_M)_++\lambda_{\rm rep}\ge\sigma_M,
       \]
       then the uniform restricted mass--degree certificate cannot satisfy
       \eqref{eq:tpc61-ledger-zero-cost-test} at the proposed exponents.
       The distinguished and signed routes remain separate questions.
 \item If one pays the complete ambient native degree, the positive gap
       \eqref{eq:tpc61-ledger-direct-native-gap} closes that proof even
       before the other endpoint losses are charged.
 \item If a literal coefficient vector satisfies
       \(A_X\zeta\ne0\) but \((A_X+M_X)\zeta=0\), the TPC--60 uniform
       shorted-Gram gate is exactly dead.
 \item If a proved necessary loss reaches \(1/400\), or the sum in
       \eqref{eq:tpc61-complete-endpoint-ledger} reaches it without a
       quantified constant-level margin, the endpoint route must stop.
\end{enumerate}

None of these failures may be repaired by averaging over shifts, replacing
the fixed \(h_0\), renormalizing the raw atoms by fiber sizes, deleting
coefficient-dependent zeros, or substituting a pointwise Rayleigh quotient
for a uniform operator norm.  Those operations change the mathematical
problem rather than close the stated gate.
